% Generated from locale/tr/content/history/set-theory/infinitesimals.tex; do not hand-edit.


\olfileid[tr]{his}{set}{infinitesimals}
\olsection{Sonsuz Küçükler ve Türev Alma}

Newton ile Leibniz kalkülüsü 17. yüzyılın sonunda (birbirlerinden bağımsız
olarak) keşfettiler. Kalkülüsün özellikle önemli bir uygulaması \emph{türev
alma} idi. Kabaca söylersek türev alma, bir fonksiyonun belirli bir noktadaki
``değişim hızı'' ya da eğimi için bir kavram vermeyi amaçlar.

Fikri canlı biçimde gösterecek bir yol şudur. Aşağıda siyahla çizilmiş
$f(x) = \nicefrac{x^2}{4} + \nicefrac{1}{2}$ fonksiyonunu ele alın:
\begin{center}
	\begin{tikzpicture}[scale=1]
	\draw[->, gray] (-1,0) -- (4.25,0) node[right] {$x$};
	\draw[->, gray] (0,-0.25) -- (0,5.25) node[above] {$f(x)$};
	
	\foreach \x/\xtext in {1/1, 2/2, 3/3, 4/4}
	\draw[shift={(\x,0)}, gray] (0pt,2pt) -- (0pt,-2pt) node[below] {$\xtext$};
	
	\foreach \y/\ytext in {1/1, 2/2, 3/3, 4/4, 5/5}
	\draw[shift={(0,\y)}, gray] (2pt,0pt) -- (-2pt,0pt) node[left] {$\ytext$};
	
	\draw[oldiagcolorC] (0.5, 9/16) -- (3.5, 57/16) -- (3.5, 9/16)--cycle;
	\draw[oldiagcolorD] (.5, 9/16) -- (2.5, 33/16) -- (2.5, 9/16) -- cycle;
	\draw[oldiagcolorE] (.5, 9/16) -- (1.5, 17/16) -- (1.5, 9/16) -- cycle;
	\draw[thick] (-1,.75) parabola bend (0,0.5) (4,4.5);
	\end{tikzpicture}
\end{center}
Fonksiyonun $c = \nicefrac{1}{2}$ noktasındaki eğimini bulmak istediğimizi
varsayalım. İşe, hipotenüsü o noktadaki eğime yaklaşan bir üçgen, örneğin
yukarıdaki kırmızı üçgeni çizerek başlarız. Üçgenimizin taban uzunluğu
$\beta$ olduğunda yüksekliği $f(\nicefrac{1}{2}+\beta) -
f(\nicefrac{1}{2})$ olur; dolayısıyla hipotenüsün eğimi şudur:
\[
\frac{f(\nicefrac{1}{2}+\beta) - f(\nicefrac{1}{2})}{\beta}.
\]
Öyleyse taban uzunluğu~$3$ olan !!{colorC} üçgenimizin eğimi tam olarak~$1$'dir.
Daha küçük bir üçgenin, taban uzunluğu~$2$ olan !!{colorD} üçgenin hipotenüsü
daha iyi bir yaklaşım verir; eğimi $\nicefrac{3}{4}$'tür. Daha da küçük,
taban uzunluğu~$1$ olan !!{colorE} üçgen daha da iyi bir yaklaşım verir;
eğimi $\nicefrac{1}{2}$'dir.

Giderek küçülen üçgenler giderek daha iyi yaklaşımlar verir. Bu nedenle şöyle
bir şey söyleyebiliriz: \emph{sonsuz küçük} taban uzunluğuna sahip bir
üçgenin hipotenüsü bize $c = \nicefrac{1}{2}$ noktasındaki eğimin kendisini
verir. Böylece $f$ fonksiyonunun $c$ noktasındaki (birinci) türevi için şu
formülü elde ederiz:
\[
{f'}(c) = \frac{f(c+\beta) - f(c)}{\beta} \text{; burada $\beta$ sonsuz küçüktür.}
\]
Newton ile Leibniz'in söyledikleri de kabaca buydu.

Ne var ki bunu söylediklerine göre onlara şunu sormalıyız: \emph{sonsuz
küçük} nedir? Ciddi bir ikilem doğar. $\beta = 0$ ise $f'$ iyi tanımlı
değildir; çünkü $0$'a bölme içerir. Ama $\beta > 0$ ise eğimin kendisini değil,
yalnızca eğime bir \emph{yaklaşımı} elde ederiz.

Bu, geçmişe bugünün ölçütlerini dayatan bir kaygı değildir. Berkeley,
Newton'ın izleyicilerini şöyle eleştirir:
\begin{quote}
İşaretlerin herhangi bir şeyi de hiçbir şeyi de gösterebileceğini kabul
ediyorum: dolayısıyla özgün $c + \beta$ gösteriminde $\beta$ ya bir artışı ya
da hiçliği gösterebilirdi. Fakat ona bunlardan hangisini göstertecekseniz o
anlamla tutarlı biçimde akıl yürütmeli, çift anlam üzerinden ilerlememelisiniz:
bunu yapmak açık bir safsata olurdu. (\citealt[\S{}XIII]{Berkeley1734},
değişkenler önceki metne uyacak biçimde değiştirilmiştir)
\end{quote}
Sonsuz küçükler kalkülüsünü Berkeley'e karşı savunmak için, ``sonsuz
küçükler''den söz etmenin yalnızca mecaz olduğu yanıtı verilebilir. \emph{Çok
küçük} bir üçgen aldığımız sürece teğete \emph{yeterince iyi} bir yaklaşım
elde edeceğimiz söylenebilir. Berkeley'in buna da bir yanıtı vardı: bu,
mühendislik için yeterince iyi olsa bile matematiğin \emph{konumunu} sarsar;
çünkü
\begin{quote}
bize \emph{in rebus mathematicis errores qu\`{a}m minimi non sunt
contemnendi} denir. [Matematik söz konusu olduğunda en küçük hatalar bile göz
ardı edilmemelidir.] \citep[\S{}IX]{Berkeley1734}
\end{quote}
İtalik bölüm, Newton'ın kendi \emph{Quadrature of Curves} (1704) eserinden
neredeyse sözcüğü sözcüğüne bir alıntıdır.

Berkeley'in felsefi itirazları son derece keskindir. Yine de kalkülüs çok büyük
başarı kazanmış bir girişimdi ve matematikçiler yanlışa düşmeden onu
kullanmayı sürdürdüler.

